\shead{Thu thập và biểu diễn số liệu}
% ---- chapter 11 helpers: bar chart, pie chart with legend, pictograph icons ----
\newcommand{\stpBars}[6]{% #1 cm per unit, #2 ymax, #3 ystep, #4 list name/height/shown, #5 y caption, #6 x caption
\begin{tikzpicture}[font=\footnotesize,y=#1cm,x=1cm,>=stealth]
\foreach \n/\v/\d [count=\i] in {#4}{\xdef\stpN{\i}}
\pgfmathsetmacro\stpW{\stpN*1.2+0.1}
\foreach \t in {0,#3,...,#2}{\draw[Line] (0,\t)--(\stpW,\t); \node[left,text=Slate] at (0,\t){$\t$};}
\foreach \n/\v/\d [count=\i] in {#4}{
  \pgfmathsetmacro\stpC{\i*1.2-0.5}
  \fill[tmBlue!65] (\stpC-0.3,0) rectangle (\stpC+0.3,\v);
  \node[above,font=\footnotesize\bfseries,text=Ink] at (\stpC,\v){$\d$};
  \node[below,text=Ink] at (\stpC,0){\n};}
\draw[->,line width=0.8pt] (0,0)--([yshift=0.55cm]0,#2);
\draw[->,line width=0.8pt] (0,0)--(\stpW+0.35,0);
\node[above,text=Slate] at ([yshift=0.55cm]0,#2){#5};
\node[right,text=Slate] at (\stpW+0.35,0){#6};
\end{tikzpicture}}
\newcommand{\stpPie}[2][1.5]{% #2 list value/text/colour/legend name (value in percent)
\begin{tikzpicture}[font=\footnotesize]
\def\stpA{90}
\foreach \v/\t/\c/\n [count=\i] in {#2}{
  \pgfmathsetmacro\stpB{\stpA-3.6*\v}
  \pgfmathsetmacro\stpM{(\stpA+\stpB)/2}
  \pgfmathsetmacro\stpR{0.64*#1}
  \fill[\c] (0,0) -- (\stpA:#1) arc[start angle=\stpA,end angle=\stpB,radius=#1] -- cycle;
  \draw[white,line width=1.2pt] (0,0) -- (\stpA:#1);
  \node[text=Ink,font=\footnotesize\bfseries] at (\stpM:\stpR) {\t};
  \fill[\c] (#1+0.7,#1-0.55*\i+0.1) rectangle ++(0.32,0.32);
  \node[anchor=west,text=Ink,font=\small] at (#1+1.1,#1-0.55*\i+0.26) {\n};
  \global\let\stpA\stpB}
\draw[Slate,line width=0.8pt] (0,0) circle (#1);
\end{tikzpicture}}
\newcommand{\stpIcons}[1]{\foreach \k in {1,...,#1}{\textcolor{tmBlue}{\faBook}\ }}
\begin{knowbox}
Để so sánh, đánh giá và đưa ra quyết định, người ta \textbf{thu thập thông tin} (số liệu) rồi \textbf{sắp xếp, xử lí} các số liệu đó. Có khi ta thu thập toàn bộ, có khi chỉ thu thập ở một nhóm có lựa chọn mang tính điển hình.\par\smallskip
Các hình thức biểu diễn số liệu thường dùng:
\begin{itemize}[leftmargin=1.4em,itemsep=3pt]
\item \textbf{bảng số liệu} -- mỗi số liệu được ghi vào một ô;
\item \textbf{biểu đồ tranh} -- mỗi hình vẽ biểu thị cùng một số lượng;
\item \textbf{biểu đồ cột} -- chiều cao mỗi cột, đọc theo vạch chia trên trục đứng, cho biết số liệu;
\item \textbf{biểu đồ hình quạt tròn} -- cả hình tròn biểu thị toàn bộ ($100\%$); mỗi hình quạt biểu thị một tỉ số phần trăm của toàn bộ. Nửa hình tròn là $50\%$, một phần tư hình tròn là $25\%$.
\end{itemize}
Với biểu đồ hình quạt: \quad \textbf{số phần $=$ tổng số $:100\times$ số phần trăm}; \quad \textbf{tổng số $=$ số phần $:$ số phần trăm $\times100$}.
\end{knowbox}
\shead{Số trung bình cộng}
\begin{knowbox}
\textbf{Số trung bình cộng} $=$ (tổng các số liệu) $:$ (số các số liệu). \qquad \textbf{Tổng} $=$ số trung bình cộng $\times$ số các số liệu.
\begin{itemize}[leftmargin=1.4em,itemsep=3pt]
\item Số trung bình cộng luôn nằm giữa số liệu bé nhất và số liệu lớn nhất.
\item Muốn tìm \emph{một số liệu chưa biết}: tính tổng từ số trung bình cộng, rồi trừ đi các số liệu đã biết.
\item Nếu một số liệu tăng thêm bao nhiêu thì tổng tăng thêm bấy nhiêu, còn số trung bình cộng tăng thêm (số đó) $:$ (số các số liệu).
\end{itemize}
\end{knowbox}
\shead{Khả năng xảy ra và đếm các kết quả}
\begin{knowbox}
Khi thực hiện lặp lại một hoạt động (tung đồng xu, gieo xúc xắc, quay vòng quay, lấy bóng mà không nhìn), ta ghi lại số lần xuất hiện của mỗi khả năng.
\begin{center}
\textbf{tỉ số mô tả số lần lặp lại của một khả năng} $=$ $\dfrac{\text{số lần khả năng đó xảy ra}}{\text{tổng số lần thực hiện}}$
\end{center}
Ví dụ: gieo xúc xắc $20$ lần, mặt $6$ chấm xuất hiện $3$ lần thì tỉ số đó là $\dfrac{3}{20}$. Nếu gieo thêm nhiều lần, tỉ số này có thể thay đổi.\par\smallskip
\textbf{Đếm các kết quả.} Liệt kê có thứ tự mọi kết quả có thể (bằng danh sách, bảng hoặc sơ đồ cây). Đồng xu có $2$ kết quả, xúc xắc có $6$ kết quả; tung một đồng xu \emph{và} gieo một xúc xắc thì có $2\times6=12$ kết quả.\par
Khi các kết quả có \emph{khả năng như nhau}, \textbf{khả năng xảy ra} của một sự kiện $=$ (số kết quả làm sự kiện xảy ra) $:$ (số tất cả các kết quả). Sự kiện luôn xảy ra là \emph{chắc chắn}; sự kiện không bao giờ xảy ra là \emph{không thể}.
\end{knowbox}
\gap
\begin{exbox}{\ExWord\ 1}
Lan ghi lại số trang sách mình đọc mỗi ngày trong một tuần:
\begin{center}
\begin{tabular}{|l|c|c|c|c|c|c|c|}
\hline
\textbf{Thứ} & Hai & Ba & Tư & Năm & Sáu & Bảy & CN\\\hline
\textbf{Số trang} & $18$ & $25$ & $20$ & $32$ & $15$ & $30$ & $28$\\\hline
\end{tabular}
\end{center}
\begin{cols2}
\pa{a} Ngày nào Lan đọc nhiều nhất? Ít nhất? & \pa{b} Trung bình mỗi ngày Lan đọc bao nhiêu trang?\\
\multicolumn{2}{@{}l@{}}{\pa{c} Có bao nhiêu ngày Lan đọc nhiều hơn số trang trung bình?}
\end{cols2}
\guide
\pa{a} Nhiều nhất vào thứ Năm ($32$ trang); ít nhất vào thứ Sáu ($15$ trang).\par
\pa{b} Số trang đọc trong tuần: $18+25+20+32+15+30+28=168$ (trang).\par Trung bình mỗi ngày: $168:7=24$ (trang).\par
\pa{c} Nhiều hơn $24$ trang: thứ Ba, thứ Năm, thứ Bảy và Chủ nhật, tức là $4$ ngày.
\dapso{\pa{a} Thứ Năm; thứ Sáu. \pa{b} $24$ trang. \pa{c} $4$ ngày.}
\end{exbox}
\gap
\begin{exbox}{\ExWord\ 2}
Biểu đồ cột dưới đây cho biết số sách năm lớp quyên góp cho thư viện trường.
\begin{center}
\stpBars{0.05}{80}{10}{5A/45/45,5B/60/60,5C/35/35,5D/70/70,5E/40/40}{(quyển)}{(lớp)}
\end{center}
\begin{cols2}
\pa{a} Cả năm lớp quyên góp được bao nhiêu quyển sách? & \pa{b} Trung bình mỗi lớp quyên góp bao nhiêu quyển?\\
\multicolumn{2}{@{}l@{}}{\pa{c} Lớp góp nhiều nhất hơn lớp góp ít nhất bao nhiêu quyển sách?}
\end{cols2}
\guide
\pa{a} $45+60+35+70+40=250$ (quyển).\par
\pa{b} $250:5=50$ (quyển).\par
\pa{c} Nhiều nhất: lớp 5D ($70$ quyển); ít nhất: lớp 5C ($35$ quyển).\par Hơn nhau: $70-35=35$ (quyển).
\dapso{\pa{a} $250$ quyển. \pa{b} $50$ quyển. \pa{c} $35$ quyển.}
\end{exbox}
\gap
\begin{exbox}{\ExWord\ 3}
Khảo sát $400$ học sinh về loại quả yêu thích nhất, kết quả được cho trong biểu đồ hình quạt tròn.
\begin{center}
\stpPie{35/35\%/tmBlue!55/Xoài,20/20\%/tmOrange!60/Chuối,25/25\%/tmGreen!50/Dưa hấu,20/?/tmPurple!40/Cam}
\end{center}
\pa{a} Có bao nhiêu phần trăm học sinh chọn cam?\par
\pa{b} Có bao nhiêu học sinh chọn xoài?\par
\pa{c} Số học sinh chọn xoài nhiều hơn số học sinh chọn chuối bao nhiêu bạn?
\guide
\pa{a} Cả hình tròn là $100\%$: $100\%-35\%-20\%-25\%=20\%$.\par
\pa{b} $400:100\times35=140$ (học sinh).\par
\pa{c} \emph{Cách 1.} Số học sinh chọn chuối: $400:100\times20=80$ (học sinh); \ $140-80=60$ (học sinh).\par
\emph{Cách 2.} Chọn xoài nhiều hơn chọn chuối $35\%-20\%=15\%$; \ $400:100\times15=60$ (học sinh).
\dapso{\pa{a} $20\%$. \pa{b} $140$ học sinh. \pa{c} $60$ học sinh.}
\end{exbox}
\gap
\begin{exbox}{\ExWord\ 4}
Minh gieo một con xúc xắc $40$ lần và ghi lại kết quả.
\begin{center}
\begin{tabular}{|l|c|c|c|c|c|c|}
\hline
\textbf{Khả năng} & $1$ chấm & $2$ chấm & $3$ chấm & $4$ chấm & $5$ chấm & $6$ chấm\\\hline
\textbf{Số lần} & $6$ & $8$ & $5$ & $7$ & $9$ & $5$\\\hline
\end{tabular}
\end{center}
Viết tỉ số mô tả số lần lặp lại của mỗi khả năng sau so với tổng số lần gieo:\par
\pa{a} mặt $6$ chấm; \quad \pa{b} mặt có số chấm chẵn; \quad \pa{c} mặt có ít nhất $5$ chấm.
\guide
\pa{a} Mặt $6$ chấm xuất hiện $5$ lần: $\dfrac{5}{40}=\dfrac{1}{8}$.\par
\pa{b} Số chấm chẵn ($2$, $4$, $6$ chấm): $8+7+5=20$ (lần); \ $\dfrac{20}{40}=\dfrac12$.\par
\pa{c} Ít nhất $5$ chấm ($5$ hoặc $6$ chấm): $9+5=14$ (lần); \ $\dfrac{14}{40}=\dfrac{7}{20}$.
\note{Điều này không có nghĩa là lần sau cũng xảy ra như vậy: gieo càng nhiều lần thì các tỉ số này càng thay đổi.}
\dapso{\pa{a} $\dfrac18$. \pa{b} $\dfrac12$. \pa{c} $\dfrac{7}{20}$.}
\end{exbox}
\gap
\begin{exbox}{\ExWord\ 5}
Hộp I có ba tấm thẻ ghi số $1$, $2$, $3$; hộp II có hai tấm thẻ ghi số $4$, $5$. Không nhìn vào hộp, Hoa lấy ở mỗi hộp một tấm thẻ rồi cộng hai số trên thẻ.\par
\pa{a} Liệt kê tất cả các kết quả có thể. \quad \pa{b} Khả năng tổng bằng $7$ là bao nhiêu?
\guide
\pa{a} Lập bảng (hàng: thẻ hộp I, cột: thẻ hộp II):
\begin{center}
\begin{tabular}{|c|c|c|}
\hline
$+$ & $4$ & $5$\\\hline
$1$ & $5$ & $6$\\\hline
$2$ & $6$ & $7$\\\hline
$3$ & $7$ & $8$\\\hline
\end{tabular}
\end{center}
Có $3\times2=6$ kết quả có khả năng như nhau: $1+4;\ 1+5;\ 2+4;\ 2+5;\ 3+4;\ 3+5$.\par
\pa{b} Tổng bằng $7$ ở $2$ kết quả ($2+5$ và $3+4$).\par Khả năng: $\dfrac26=\dfrac13$.
\dapso{\pa{a} $6$ kết quả. \pa{b} $\dfrac13$.}
\end{exbox}
\dividerbar{Luyện tập}
\needspace{5cm}
\begin{probbox}{1}{\relax}
Nhiệt độ lúc giữa trưa ở một thị trấn trong năm ngày như sau:
\begin{center}
\begin{tabular}{|l|c|c|c|c|c|}
\hline
\textbf{Thứ} & Hai & Ba & Tư & Năm & Sáu\\\hline
\textbf{Nhiệt độ ($^\circ$C)} & $31$ & $29$ & $33$ & $30$ & $32$\\\hline
\end{tabular}
\end{center}
\pa{a} Ngày nào nóng nhất, ngày nào mát nhất? \quad \pa{b} Tính nhiệt độ trung bình lúc giữa trưa.
\end{probbox}
\probsol{
\pa{a} Nóng nhất: thứ Tư ($33^\circ$C); mát nhất: thứ Ba ($29^\circ$C).\par
\pa{b} $(31+29+33+30+32):5=155:5=31\ (^\circ\text{C})$.
\dapso{\pa{a} Thứ Tư; thứ Ba. \pa{b} $31^\circ$C.}}
\gap
\begin{probbox}{2}{\relax}
Điểm sáu bài kiểm tra Toán của Nam là $7;\ 9;\ 8;\ 10;\ 6;\ 8$. Tính điểm trung bình của Nam.
\end{probbox}
\probsol{
Tổng điểm: $7+9+8+10+6+8=48$.\par Điểm trung bình: $48:6=8$.
\dapso{$8$ điểm.}}
\gap
\begin{probbox}{3}{\relax}
Bốn bạn cao $135$ cm, $142$ cm, $138$ cm và $145$ cm. Tính chiều cao trung bình của bốn bạn.
\end{probbox}
\probsol{
$(135+142+138+145):4=560:4=140$ (cm).
\dapso{$140$ cm.}}
\gap
\needspace{5cm}
\begin{probbox}{4}{\relax}
Biểu đồ tranh dưới đây cho biết số sách bốn tổ đã đọc trong tháng Năm.
\begin{center}
\begin{tabular}{|l|l|}
\hline
Tổ 1 & \stpIcons{5}\\\hline
Tổ 2 & \stpIcons{3}\\\hline
Tổ 3 & \stpIcons{6}\\\hline
Tổ 4 & \stpIcons{4}\\\hline
\end{tabular}\par\smallskip
Mỗi \textcolor{tmBlue}{\faBook} chỉ $4$ quyển sách.
\end{center}
\pa{a} Cả bốn tổ đọc được bao nhiêu quyển sách? \quad \pa{b} Trung bình mỗi tổ đọc bao nhiêu quyển?
\end{probbox}
\probsol{
Số hình vẽ: $5+3+6+4=18$ (hình).\par
\pa{a} Số sách: $4\times18=72$ (quyển).\par
\pa{b} Trung bình: $72:4=18$ (quyển).
\dapso{\pa{a} $72$ quyển. \pa{b} $18$ quyển.}}
\gap
\begin{probbox}{5}{\relax}
Một cửa hàng bán được $24$ chiếc xe đạp trong tháng Một, $30$ chiếc trong tháng Hai, $27$ chiếc trong tháng Ba và $39$ chiếc trong tháng Tư.\par
\pa{a} Bốn tháng cửa hàng bán được bao nhiêu chiếc xe đạp? \quad \pa{b} Trung bình mỗi tháng bán được bao nhiêu chiếc?\par
\pa{c} Có mấy tháng cửa hàng bán được nhiều hơn mức trung bình?
\end{probbox}
\probsol{
\pa{a} $24+30+27+39=120$ (chiếc).\par
\pa{b} $120:4=30$ (chiếc).\par
\pa{c} Chỉ có tháng Tư ($39>30$; tháng Hai bán $30$ chiếc, bằng mức trung bình chứ không nhiều hơn): $1$ tháng.
\dapso{\pa{a} $120$ chiếc. \pa{b} $30$ chiếc. \pa{c} $1$ tháng.}}
\gap
\needspace{8.5cm}
\begin{probbox}{6}{\relax}
Biểu đồ cho biết số học sinh mỗi khối của một trường tiểu học.
\begin{center}
\stpBars{0.025}{150}{25}{Khối 1/120/120,Khối 2/135/135,Khối 3/110/110,Khối 4/140/140,Khối 5/125/125}{(học sinh)}{}
\end{center}
\pa{a} Khối nào có nhiều học sinh nhất? Ít học sinh nhất?\par
\pa{b} Trường có tất cả bao nhiêu học sinh?\par
\pa{c} Khối 4 nhiều hơn khối 3 bao nhiêu học sinh?
\end{probbox}
\probsol{
\pa{a} Nhiều nhất: khối 4 ($140$ học sinh); ít nhất: khối 3 ($110$ học sinh).\par
\pa{b} $120+135+110+140+125=630$ (học sinh).\par
\pa{c} $140-110=30$ (học sinh).
\dapso{\pa{a} Khối 4; khối 3. \pa{b} $630$ học sinh. \pa{c} $30$ học sinh.}}
\gap
\needspace{8.5cm}
\begin{probbox}{7}{\relax}
Biểu đồ cho biết lượng mưa ở một thị trấn từ tháng Năm đến tháng Tám.
\begin{center}
\stpBars{0.014}{250}{50}{Tháng 5/80/80,Tháng 6/150/150,Tháng 7/210/210,Tháng 8/120/120}{(mm)}{}
\end{center}
Tính tổng lượng mưa và lượng mưa trung bình mỗi tháng.
\end{probbox}
\probsol{
Tổng lượng mưa: $80+150+210+120=560$ (mm).\par Trung bình mỗi tháng: $560:4=140$ (mm).
\dapso{$560$ mm; $140$ mm.}}
\gap
\needspace{7.5cm}
\begin{probbox}{8}{\relax}
Lớp 5B có $40$ học sinh. Biểu đồ hình quạt tròn cho biết cách các bạn đến trường.
\begin{center}
\stpPie{30/30\%/tmBlue!55/Đi bộ,45/45\%/tmOrange!60/Đi xe đạp,25/25\%/tmGreen!50/Bố mẹ đưa bằng xe máy}
\end{center}
\pa{a} Có bao nhiêu học sinh đi bộ đến trường?\par
\pa{b} Số học sinh đi xe đạp nhiều hơn số học sinh được đưa bằng xe máy bao nhiêu bạn?
\end{probbox}
\probsol{
\pa{a} $40:100\times30=12$ (học sinh).\par
\pa{b} Đi xe đạp: $40:100\times45=18$ (học sinh); đi xe máy: $40:100\times25=10$ (học sinh).\par Nhiều hơn: $18-10=8$ (học sinh).
\dapso{\pa{a} $12$ học sinh. \pa{b} $8$ học sinh.}}
\gap
\needspace{7.5cm}
\begin{probbox}{9}{\relax}
$200$ học sinh chọn môn thể thao yêu thích nhất của mình.
\begin{center}
\stpPie{40/40\%/tmBlue!55/Bóng đá,25/25\%/tmOrange!60/Cầu lông,20/20\%/tmGreen!50/Bơi,15/15\%/tmPurple!40/Cờ vua}
\end{center}
\pa{a} Có bao nhiêu học sinh chọn bóng đá? \quad \pa{b} chọn cờ vua?\par
\pa{c} Có bao nhiêu học sinh chọn bóng đá hoặc bơi?
\end{probbox}
\probsol{
\pa{a} $200:100\times40=80$ (học sinh).\par
\pa{b} $200:100\times15=30$ (học sinh).\par
\pa{c} $40\%+20\%=60\%$; \ $200:100\times60=120$ (học sinh).
\dapso{\pa{a} $80$. \pa{b} $30$. \pa{c} $120$ học sinh.}}
\gap
\needspace{7.5cm}
\begin{probbox}{10}{\relax}
$300$ học sinh, mỗi bạn tham gia một câu lạc bộ sau giờ học.
\begin{center}
\stpPie{30/30\%/tmBlue!55/Mĩ thuật,25/25\%/tmOrange!60/Âm nhạc,25/?/tmGreen!50/Khoa học,20/20\%/tmPurple!40/Thể thao}
\end{center}
Có bao nhiêu phần trăm học sinh tham gia câu lạc bộ Khoa học và đó là bao nhiêu học sinh?
\end{probbox}
\probsol{
Khoa học: $100\%-30\%-25\%-20\%=25\%$.\par Số học sinh: $300:100\times25=75$ (học sinh).
\dapso{$25\%$; $75$ học sinh.}}
\gap
\needspace{7.5cm}
\begin{probbox}{11}{\relax}
Một gia đình chi tiêu $12\,000\,000$ đồng mỗi tháng như biểu đồ dưới đây.
\begin{center}
\stpPie{50/50\%/tmBlue!55/Ăn uống,25/25\%/tmOrange!60/Học tập,15/15\%/tmGreen!50/Tiết kiệm,10/10\%/tmPurple!40/Việc khác}
\end{center}
Gia đình chi cho ăn uống bao nhiêu tiền? Để tiết kiệm bao nhiêu tiền?
\end{probbox}
\probsol{
Ăn uống chiếm một nửa: $12\,000\,000:2=6\,000\,000$ (đồng).\par
Tiết kiệm: $12\,000\,000:100\times15=1\,800\,000$ (đồng).
\dapso{Ăn uống $6\,000\,000$ đồng; tiết kiệm $1\,800\,000$ đồng.}}
\gap
\begin{probbox}{12}{\relax}
Trong một cuộc khảo sát, có $36$ học sinh chọn xoài là loại quả yêu thích nhất. Số học sinh này chiếm $30\%$ số học sinh được khảo sát. Hỏi có bao nhiêu học sinh được khảo sát?
\end{probbox}
\probsol{
$30\%$ số học sinh là $36$ học sinh, nên toàn bộ ($100\%$) số học sinh là: $36:30\times100=120$ (học sinh).
\dapso{$120$ học sinh.}}
\gap
\begin{probbox}{13}{\relax}
Trung bình cộng của năm số là $18$. Bốn trong năm số đó là $15;\ 20;\ 17;\ 22$. Tìm số thứ năm.
\end{probbox}
\probsol{
Tổng năm số: $18\times5=90$.\par Tổng bốn số: $15+20+17+22=74$.\par Số thứ năm: $90-74=16$.
\dapso{$16$.}}
\gap
\begin{probbox}{14}{\relax}
Điểm trung bình ba bài kiểm tra của An là $8$. Sau bài kiểm tra thứ tư, điểm trung bình của An là $8{,}5$. Hỏi bài thứ tư An được mấy điểm?
\end{probbox}
\probsol{
Tổng điểm ba bài: $8\times3=24$.\par Tổng điểm bốn bài: $8{,}5\times4=34$.\par Bài thứ tư: $34-24=10$ (điểm).
\dapso{$10$ điểm.}}
\gap
\begin{probbox}{15}{\relax}
Ba bao gạo nặng trung bình $25$ kg. Người ta để thêm một bao thứ tư nặng $33$ kg. Hỏi trung bình mỗi bao trong bốn bao nặng bao nhiêu ki-lô-gam?
\end{probbox}
\probsol{
Ba bao nặng: $25\times3=75$ (kg).\par Bốn bao nặng: $75+33=108$ (kg).\par Trung bình: $108:4=27$ (kg).
\dapso{$27$ kg.}}
\gap
\begin{probbox}{16}{\relax}
Tú tung một đồng xu $50$ lần, có $28$ lần xuất hiện mặt sấp. Viết tỉ số mô tả số lần xuất hiện mặt sấp, mặt ngửa so với tổng số lần tung.
\end{probbox}
\probsol{
Mặt sấp: $\dfrac{28}{50}=\dfrac{14}{25}$.\par Mặt ngửa: $50-28=22$ (lần), $\dfrac{22}{50}=\dfrac{11}{25}$.
\dapso{Mặt sấp $\dfrac{14}{25}$; mặt ngửa $\dfrac{11}{25}$.}}
\gap
\begin{probbox}{17}{\relax}
Một vòng quay có các phần màu đỏ, xanh và vàng được quay $60$ lần: kim dừng ở màu đỏ $21$ lần, màu xanh $15$ lần, màu vàng $24$ lần.\par
\pa{a} Viết tỉ số cho mỗi màu. \quad \pa{b} Viết tỉ số cho khả năng ``không phải màu đỏ''.
\end{probbox}
\probsol{
\pa{a} Đỏ: $\dfrac{21}{60}=\dfrac{7}{20}$; \ xanh: $\dfrac{15}{60}=\dfrac14$; \ vàng: $\dfrac{24}{60}=\dfrac25$.\par
\pa{b} Không phải màu đỏ: $15+24=39$ (lần); \ $\dfrac{39}{60}=\dfrac{13}{20}$.
\dapso{\pa{a} $\dfrac{7}{20}$; $\dfrac14$; $\dfrac25$. \pa{b} $\dfrac{13}{20}$.}}
\gap
\needspace{5cm}
\begin{probbox}{18}{\relax}
Gieo một con xúc xắc $30$ lần, được kết quả:
\begin{center}
\begin{tabular}{|l|c|c|c|c|c|c|}
\hline
\textbf{Khả năng} & $1$ chấm & $2$ chấm & $3$ chấm & $4$ chấm & $5$ chấm & $6$ chấm\\\hline
\textbf{Số lần} & $4$ & $6$ & $5$ & $3$ & $7$ & $5$\\\hline
\end{tabular}
\end{center}
Viết tỉ số cho khả năng: \pa{a} mặt có số chấm lẻ; \quad \pa{b} mặt có nhiều hơn $4$ chấm.
\end{probbox}
\probsol{
\pa{a} Số chấm lẻ ($1$, $3$, $5$): $4+5+7=16$ (lần); \ $\dfrac{16}{30}=\dfrac{8}{15}$.\par
\pa{b} Nhiều hơn $4$ chấm ($5$, $6$): $7+5=12$ (lần); \ $\dfrac{12}{30}=\dfrac25$.
\dapso{\pa{a} $\dfrac{8}{15}$. \pa{b} $\dfrac25$.}}
\gap
\begin{probbox}{19}{\relax}
Bình gieo một con xúc xắc nhiều lần. Mặt $6$ chấm xuất hiện $9$ lần và tỉ số mô tả số lần xuất hiện mặt $6$ chấm so với tổng số lần gieo là $\dfrac{3}{20}$. Hỏi Bình đã gieo xúc xắc bao nhiêu lần? Có bao nhiêu lần không ra mặt $6$ chấm?
\end{probbox}
\probsol{
$9$ lần ứng với $3$ phần trong $20$ phần bằng nhau, nên một phần là $9:3=3$ (lần).\par Số lần gieo: $3\times20=60$ (lần).\par Không ra mặt $6$ chấm: $60-9=51$ (lần).
\dapso{$60$ lần; $51$ lần.}}
\gap
\begin{probbox}{20}{\relax}
Một túi đựng các quả bóng đỏ và xanh. Mai lấy ra một quả bóng mà không nhìn, ghi lại màu rồi trả bóng vào túi. Mai làm như vậy $25$ lần và tỉ số cho màu đỏ là $\dfrac25$. Hỏi Mai lấy được bóng đỏ bao nhiêu lần, bóng xanh bao nhiêu lần?
\end{probbox}
\probsol{
Bóng đỏ: $25\times\dfrac25=10$ (lần).\par Bóng xanh: $25-10=15$ (lần).
\dapso{Bóng đỏ $10$ lần, bóng xanh $15$ lần.}}
\gap
\begin{probbox}{21}{\relax}
Tung cùng lúc hai đồng xu.\par
\pa{a} Liệt kê tất cả các kết quả có thể. \quad \pa{b} Có bao nhiêu kết quả gồm một mặt sấp và một mặt ngửa?
\end{probbox}
\probsol{
Kí hiệu S là mặt sấp, N là mặt ngửa (đồng xu thứ nhất viết trước).\par
\pa{a} SS, SN, NS, NN: $4$ kết quả.\par
\pa{b} SN và NS: $2$ kết quả.
\dapso{\pa{a} $4$ kết quả. \pa{b} $2$ kết quả.}}
\gap
\begin{probbox}{22}{\relax}
Căng tin trường có $3$ món chính (cơm, mì, bánh mì) và $2$ loại đồ uống (sữa, nước cam). Mỗi suất trưa gồm một món chính và một đồ uống. Hỏi có bao nhiêu suất trưa khác nhau? Hãy liệt kê.
\end{probbox}
\probsol{
Mỗi món chính đi với $2$ loại đồ uống: $3\times2=6$ (suất).\par
Cơm -- sữa, cơm -- nước cam, mì -- sữa, mì -- nước cam, bánh mì -- sữa, bánh mì -- nước cam.
\dapso{$6$ suất.}}
\gap
\begin{probbox}{23}{\relax}
Từ các chữ số $1;\ 4;\ 7$ có thể viết được bao nhiêu số có hai chữ số khác nhau? Trong đó có bao nhiêu số chẵn?
\end{probbox}
\probsol{
Các số: $14;\ 17;\ 41;\ 47;\ 71;\ 74$: $6$ số ($3$ cách chọn chữ số hàng chục, rồi $2$ cách chọn chữ số hàng đơn vị).\par
Số chẵn có tận cùng là $4$: $14$ và $74$, tức là $2$ số.
\dapso{$6$ số; $2$ số chẵn.}}
\gap
\begin{probbox}{24}{\relax}
Một túi có $3$ quả bóng đỏ, $5$ quả bóng xanh và $2$ quả bóng vàng có cùng kích thước. Lấy ra một quả bóng mà không nhìn.\par
\pa{a} Màu nào có khả năng được lấy ra nhiều nhất? \quad \pa{b} Khả năng lấy được bóng xanh là bao nhiêu?\par
\pa{c} Lấy được bóng màu lục là chắc chắn, có thể hay không thể?
\end{probbox}
\probsol{
\pa{a} Màu xanh, vì bóng xanh nhiều nhất ($5$ quả).\par
\pa{b} Có $10$ quả, trong đó $5$ quả xanh: $\dfrac{5}{10}=\dfrac12$.\par
\pa{c} Túi không có bóng lục: không thể.
\dapso{\pa{a} Màu xanh. \pa{b} $\dfrac12$. \pa{c} Không thể.}}
\gap
\begin{probbox}{25}{\relax}
Một đội bóng ghi được $2;\ 0;\ 3;\ 1;\ 4;\ 2;\ 1;\ 3$ bàn thắng trong tám trận. Tính số bàn thắng trung bình mỗi trận. Có bao nhiêu trận đội ghi nhiều bàn hơn mức trung bình?
\end{probbox}
\probsol{
Số bàn thắng: $2+0+3+1+4+2+1+3=16$ (bàn).\par Trung bình: $16:8=2$ (bàn).\par Nhiều hơn $2$ bàn: các trận ghi $3;\ 4;\ 3$ bàn, tức là $3$ trận.
\dapso{$2$ bàn; $3$ trận.}}
\gap
\needspace{8.5cm}
\begin{probbox}{26}{\relax}
Biểu đồ cho biết sản lượng lúa của một trang trại (đơn vị: tấn).
\begin{center}
\stpBars{0.06}{60}{10}{2021/42/42,2022/45/45,2023/48/48,2024/50/50,2025/55/55}{(tấn)}{(năm)}
\end{center}
\pa{a} Tính sản lượng lúa trung bình mỗi năm. \quad \pa{b} Năm 2025 thu hoạch nhiều hơn năm 2021 bao nhiêu tấn?
\end{probbox}
\probsol{
\pa{a} $(42+45+48+50+55):5=240:5=48$ (tấn).\par
\pa{b} $55-42=13$ (tấn).
\dapso{\pa{a} $48$ tấn. \pa{b} $13$ tấn.}}
\gap
\needspace{7.5cm}
\begin{probbox}{27}{\relax}
Một thư viện có $1\,200$ quyển sách.
\begin{center}
\stpPie{35/35\%/tmBlue!55/Truyện tranh,25/25\%/tmOrange!60/Khoa học,30/30\%/tmGreen!50/Truyện chữ,10/10\%/tmPurple!40/Loại khác}
\end{center}
\pa{a} Có bao nhiêu quyển sách khoa học? \quad \pa{b} Truyện chữ nhiều hơn sách loại khác bao nhiêu quyển?
\end{probbox}
\probsol{
\pa{a} $1\,200:100\times25=300$ (quyển).\par
\pa{b} Truyện chữ nhiều hơn $30\%-10\%=20\%$: $1\,200:100\times20=240$ (quyển).
\dapso{\pa{a} $300$ quyển. \pa{b} $240$ quyển.}}
\gap
\needspace{7.5cm}
\begin{probbox}{28}{\relax}
Hoa sử dụng $160\,000$ đồng tiền tiêu vặt như biểu đồ dưới đây.
\begin{center}
\stpPie{50/50\%/tmBlue!55/Mua sách,25/25\%/tmOrange!60/Tiết kiệm,15/15\%/tmGreen!50/Ăn vặt,10/10\%/tmPurple!40/Mua quà}
\end{center}
\pa{a} Hoa tiết kiệm được bao nhiêu tiền? \quad \pa{b} Tiền mua sách nhiều hơn tiền ăn vặt bao nhiêu?
\end{probbox}
\probsol{
\pa{a} Một phần tư số tiền: $160\,000:4=40\,000$ (đồng).\par
\pa{b} Mua sách: $80\,000$ đồng; ăn vặt: $160\,000:100\times15=24\,000$ (đồng).\par Nhiều hơn: $80\,000-24\,000=56\,000$ (đồng).
\dapso{\pa{a} $40\,000$ đồng. \pa{b} $56\,000$ đồng.}}
\gap
\begin{probbox}{29}{\relax}
Năm bạn nữ nặng trung bình $32$ kg, ba bạn nam nặng trung bình $36$ kg. Tính cân nặng trung bình của cả tám bạn.
\end{probbox}
\probsol{
Năm bạn nữ: $32\times5=160$ (kg).\par Ba bạn nam: $36\times3=108$ (kg).\par Trung bình: $(160+108):8=268:8=33{,}5$ (kg).
\dapso{$33{,}5$ kg.}}
\gap
\begin{probbox}{30}{\relax}
Lan tung một đồng xu nhiều lần. Có $14$ lần xuất hiện mặt ngửa và tỉ số cho mặt sấp là $\dfrac35$. Hỏi Lan đã tung đồng xu bao nhiêu lần và có bao nhiêu lần ra mặt sấp?
\end{probbox}
\probsol{
Tỉ số cho mặt ngửa: $1-\dfrac35=\dfrac25$.\par Vậy $14$ lần ứng với $2$ phần trong $5$ phần bằng nhau: một phần là $7$ lần.\par Số lần tung: $7\times5=35$ (lần); mặt sấp: $35-14=21$ (lần).
\dapso{$35$ lần; $21$ lần mặt sấp.}}
\dividerbar{Thử thách -- Nâng cao}
\begin{probbox}{31}{\stars{1}}
Trung bình cộng của sáu số là $15$. Bỏ đi một số thì trung bình cộng của năm số còn lại là $14$. Số bị bỏ đi là số nào?
\end{probbox}
\probsol{
Tổng sáu số: $15\times6=90$.\par Tổng năm số: $14\times5=70$.\par Số bị bỏ đi: $90-70=20$.
\dapso{$20$.}}
\gap
\begin{probbox}{32}{\stars{1}}
Tuổi trung bình của $10$ cầu thủ trong một đội là $30$. Nếu tính cả huấn luyện viên thì tuổi trung bình là $31$. Hỏi huấn luyện viên bao nhiêu tuổi?
\end{probbox}
\probsol{
Tổng số tuổi của các cầu thủ: $30\times10=300$ (tuổi).\par Tính cả huấn luyện viên ($11$ người): $31\times11=341$ (tuổi).\par Huấn luyện viên: $341-300=41$ (tuổi).
\dapso{$41$ tuổi.}}
\gap
\begin{probbox}{33}{\stars{1}}
Điểm trung bình bốn bài kiểm tra của Hùng là $7{,}5$. Bài thứ năm Hùng phải được mấy điểm để điểm trung bình là $8$?
\end{probbox}
\probsol{
Tổng điểm bốn bài: $7{,}5\times4=30$.\par Tổng điểm năm bài: $8\times5=40$.\par Bài thứ năm: $40-30=10$ (điểm).
\dapso{$10$ điểm.}}
\gap
\begin{probbox}{34}{\stars{1}}
Trung bình cộng của ba số là $35$. Số thứ nhất gấp đôi số thứ hai, số thứ hai gấp đôi số thứ ba. Tìm ba số đó.
\end{probbox}
\probsol{
Tổng ba số: $35\times3=105$.\par Coi số thứ ba là $1$ phần thì số thứ hai là $2$ phần, số thứ nhất là $4$ phần: tất cả $7$ phần.\par
Một phần: $105:7=15$.\par Ba số là: $60;\ 30;\ 15$.
\dapso{$60;\ 30;\ 15$.}}
\gap
\begin{probbox}{35}{\stars{1}}
Trung bình cộng của năm số chẵn liên tiếp là $36$. Tìm số lớn nhất trong năm số đó.
\end{probbox}
\probsol{
Với các số cách đều, trung bình cộng bằng số ở giữa: $36$.\par Năm số đó là $32;\ 34;\ 36;\ 38;\ 40$.
\dapso{$40$.}}
\gap
\begin{probbox}{36}{\stars{1}}
Điểm trung bình của bốn bài được tính ra là $8{,}25$, nhưng có một điểm bị chép nhầm thành $6$ thay vì $9$. Tính điểm trung bình đúng.
\end{probbox}
\probsol{
Tổng điểm đã tính: $8{,}25\times4=33$.\par Tổng điểm đúng: $33-6+9=36$.\par Điểm trung bình đúng: $36:4=9$.
\dapso{$9$ điểm.}}
\gap
\needspace{8.5cm}
\begin{probbox}{37}{\stars{1}}
Biểu đồ cho biết số vé bán được ở một bảo tàng từ thứ Hai đến thứ Sáu, trong đó cột của thứ Tư bị mất. Trung bình mỗi ngày bán được $160$ vé. Hỏi thứ Tư bán được bao nhiêu vé?
\begin{center}
\stpBars{0.016}{250}{50}{Thứ Hai/120/120,Thứ Ba/150/150,Thứ Tư/0/?,Thứ Năm/180/180,Thứ Sáu/210/210}{(vé)}{}
\end{center}
\end{probbox}
\probsol{
Năm ngày bán được: $160\times5=800$ (vé).\par Bốn ngày còn lại: $120+150+180+210=660$ (vé).\par Thứ Tư: $800-660=140$ (vé).
\dapso{$140$ vé.}}
\gap
\needspace{7.5cm}
\begin{probbox}{38}{\stars{1}}
Mỗi học sinh lớp 5A chọn một môn thể thao. Có $12$ học sinh chọn bóng rổ.
\begin{center}
\stpPie{40/40\%/tmBlue!55/Bóng đá,30/30\%/tmOrange!60/Bóng rổ,20/20\%/tmGreen!50/Bơi,10/10\%/tmPurple!40/Cờ vua}
\end{center}
\pa{a} Lớp 5A có bao nhiêu học sinh? \quad \pa{b} Có bao nhiêu bạn chọn bóng đá? Chọn cờ vua?
\end{probbox}
\probsol{
\pa{a} $12$ học sinh chiếm $30\%$: $12:30\times100=40$ (học sinh).\par
\pa{b} Bóng đá: $40:100\times40=16$ (bạn); \ cờ vua: $40:100\times10=4$ (bạn).
\dapso{\pa{a} $40$ học sinh. \pa{b} $16$ bạn; $4$ bạn.}}
\gap
\begin{probbox}{39}{\stars{2}}
$160$ học sinh chọn con vật nuôi yêu thích: chó $40\%$, thỏ $15\%$, số còn lại chọn mèo hoặc cá. Số bạn chọn mèo gấp đôi số bạn chọn cá. Hỏi có bao nhiêu bạn chọn mèo?
\end{probbox}
\probsol{
Mèo và cá: $100\%-40\%-15\%=45\%$.\par Cá $1$ phần, mèo $2$ phần: một phần là $45\%:3=15\%$.\par Mèo: $30\%$, tức là $160:100\times30=48$ (bạn).
\dapso{$48$ bạn.}}
\gap
\begin{probbox}{40}{\stars{1}}
Lớp 5A ($40$ học sinh) có $35\%$ số bạn thích đọc sách nhất. Lớp 5B ($36$ học sinh) có $50\%$ số bạn thích đọc sách nhất. Lớp nào có nhiều bạn thích đọc sách nhất hơn và nhiều hơn bao nhiêu bạn?
\end{probbox}
\probsol{
Lớp 5A: $40:100\times35=14$ (bạn).\par Lớp 5B: $36:2=18$ (bạn).\par Lớp 5B nhiều hơn $18-14=4$ (bạn).
\dapso{Lớp 5B, nhiều hơn $4$ bạn.}}
\gap
\begin{probbox}{41}{\stars{1}}
Từ các chữ số $0;\ 3;\ 5;\ 8$ có thể viết được bao nhiêu số có hai chữ số khác nhau? Trong đó có bao nhiêu số chia hết cho $5$?
\end{probbox}
\probsol{
Chữ số hàng chục: $3$ cách chọn ($3$, $5$, $8$; không chọn $0$); chữ số hàng đơn vị: $3$ cách chọn (ba chữ số còn lại).\par Số các số: $3\times3=9$ (số).\par
Chia hết cho $5$ là các số tận cùng bằng $0$ hoặc $5$: $30;\ 50;\ 80;\ 35;\ 85$: $5$ số.
\dapso{$9$ số; $5$ số.}}
\gap
\begin{probbox}{42}{\stars{1}}
Sáu bạn gặp nhau, cứ hai bạn thì bắt tay nhau một lần. Hỏi có tất cả bao nhiêu cái bắt tay?
\end{probbox}
\probsol{
Mỗi bạn bắt tay $5$ bạn: $6\times5=30$, nhưng mỗi cái bắt tay được tính hai lần.\par Số cái bắt tay: $30:2=15$.
\dapso{$15$ cái bắt tay.}}
\gap
\begin{probbox}{43}{\stars{2}}
Gieo cùng lúc hai con xúc xắc rồi cộng số chấm.\par
\pa{a} Có tất cả bao nhiêu kết quả? \quad \pa{b} Có bao nhiêu kết quả có tổng bằng $7$? Khả năng tổng bằng $7$ là bao nhiêu?\par
\pa{c} Khả năng tổng bằng $12$ là bao nhiêu?
\end{probbox}
\probsol{
\pa{a} Mỗi mặt trong $6$ mặt của xúc xắc thứ nhất đi với $6$ mặt của xúc xắc thứ hai: $6\times6=36$ (kết quả).\par
\pa{b} Tổng bằng $7$: $1+6;\ 2+5;\ 3+4;\ 4+3;\ 5+2;\ 6+1$: $6$ kết quả. Khả năng: $\dfrac{6}{36}=\dfrac16$.\par
\pa{c} Chỉ có $6+6$: khả năng $\dfrac{1}{36}$.
\dapso{\pa{a} $36$. \pa{b} $6$; $\dfrac16$. \pa{c} $\dfrac{1}{36}$.}}
\gap
\begin{probbox}{44}{\stars{1}}
Một hộp có $4$ quả bóng đỏ và $6$ quả bóng xanh. \pa{a} Khả năng lấy được bóng đỏ (không nhìn) là bao nhiêu?\par
\pa{b} Phải bỏ thêm bao nhiêu quả bóng đỏ vào hộp để khả năng lấy được bóng đỏ là $\dfrac12$?
\end{probbox}
\probsol{
\pa{a} $\dfrac{4}{10}=\dfrac25$.\par
\pa{b} Khả năng là $\dfrac12$ khi số bóng đỏ bằng số bóng xanh: $6$ quả đỏ.\par Phải thêm $6-4=2$ (quả bóng đỏ).
\dapso{\pa{a} $\dfrac25$. \pa{b} $2$ quả bóng đỏ.}}
\gap
\begin{probbox}{45}{\stars{1}}
Một túi có $20$ quả bóng màu đỏ, xanh và lục. Khả năng lấy được bóng đỏ là $\dfrac14$, khả năng lấy được bóng xanh là $\dfrac25$. Hỏi mỗi màu có bao nhiêu quả bóng?
\end{probbox}
\probsol{
Bóng đỏ: $20\times\dfrac14=5$ (quả).\par Bóng xanh: $20\times\dfrac25=8$ (quả).\par Bóng lục: $20-5-8=7$ (quả).
\dapso{$5$ quả đỏ, $8$ quả xanh, $7$ quả lục.}}
\gap
\needspace{7.5cm}
\begin{probbox}{46}{\stars{1}}
Một vòng quay được chia thành $8$ phần bằng nhau: $3$ phần màu đỏ, $4$ phần màu xanh và $1$ phần màu vàng.
\begin{center}
\begin{tikzpicture}[font=\footnotesize]
\foreach \c [count=\i from 0] in {tmPink!45,tmBlue!50,tmPink!45,tmBlue!50,StarGold,tmBlue!50,tmPink!45,tmBlue!50}{
  \fill[\c] (0,0) -- (90-45*\i:1.3) arc[start angle=90-45*\i,end angle=45-45*\i,radius=1.3] -- cycle;
  \draw[white,line width=1.2pt] (0,0) -- (90-45*\i:1.3);}
\draw[Slate,line width=0.8pt] (0,0) circle (1.3);
\draw[-{Stealth},line width=1.4pt,Ink] (0,0) -- (25:1.0);
\fill[Ink] (0,0) circle (2.5pt);
\foreach \c/\n [count=\i] in {tmPink!45/Đỏ,tmBlue!50/Xanh,StarGold/Vàng}{
  \fill[\c] (2.0,1.0-0.55*\i) rectangle ++(0.32,0.32);
  \node[anchor=west,font=\small] at (2.4,1.16-0.55*\i) {\n};}
\end{tikzpicture}
\end{center}
\pa{a} Khả năng kim dừng ở màu xanh là bao nhiêu? \quad \pa{b} Nếu quay $80$ lần, em dự đoán kim dừng ở màu xanh khoảng bao nhiêu lần?
\end{probbox}
\probsol{
\pa{a} $4$ trong $8$ phần bằng nhau là màu xanh: $\dfrac48=\dfrac12$.\par
\pa{b} Khoảng một nửa số lần quay: $80\times\dfrac12=40$ (lần).
\dapso{\pa{a} $\dfrac12$. \pa{b} Khoảng $40$ lần.}}
\gap
\begin{probbox}{47}{\stars{1}}
Hiện nay tuổi trung bình của một gia đình bốn người là $24$ tuổi. Hỏi $4$ năm trước tuổi trung bình của gia đình đó là bao nhiêu (khi đó cả bốn người đều đã ra đời)?
\end{probbox}
\probsol{
Tổng số tuổi hiện nay: $24\times4=96$ (tuổi).\par $4$ năm trước mỗi người kém hiện nay $4$ tuổi: $96-4\times4=80$ (tuổi).\par Tuổi trung bình: $80:4=20$ (tuổi).
\dapso{$20$ tuổi.}}
\gap
\begin{probbox}{48}{\stars{2}}
Trong một bài kiểm tra, điểm trung bình của $32$ học sinh một lớp là $8$. Điểm trung bình của $20$ bạn nữ là $8{,}3$. Tính điểm trung bình của các bạn nam.
\end{probbox}
\probsol{
Tổng điểm cả lớp: $8\times32=256$.\par Tổng điểm các bạn nữ: $8{,}3\times20=166$.\par Các bạn nam: $32-20=12$ (bạn), tổng điểm $256-166=90$.\par Điểm trung bình: $90:12=7{,}5$.
\dapso{$7{,}5$ điểm.}}
\gap
\begin{probbox}{49}{\stars{1}}
Một hộp có $12$ quả bóng, gồm bóng đỏ và bóng trắng. Nam lấy một quả bóng, ghi màu rồi trả lại hộp, làm như vậy $60$ lần; có $45$ lần lấy được bóng đỏ. Hãy ước lượng trong hộp có khoảng bao nhiêu quả bóng đỏ.
\end{probbox}
\probsol{
Tỉ số cho bóng đỏ: $\dfrac{45}{60}=\dfrac34$.\par Vậy khoảng $\dfrac34$ số bóng là bóng đỏ: $12\times\dfrac34=9$ (quả).
\dapso{Khoảng $9$ quả bóng đỏ.}}
\gap
\begin{probbox}{50}{\stars{1}}
An, Bình và Chi đứng thành một hàng để chụp ảnh.\par
\pa{a} Có bao nhiêu cách đứng khác nhau? \quad \pa{b} Có bao nhiêu cách mà An đứng ở giữa?
\end{probbox}
\probsol{
\pa{a} Vị trí thứ nhất: $3$ cách chọn, thứ hai: $2$ cách, thứ ba: $1$ cách; \ $3\times2\times1=6$ (cách):\par
ABC, ACB, BAC, BCA, CAB, CBA (A là An, B là Bình, C là Chi).\par
\pa{b} An đứng giữa: BAC và CAB, tức là $2$ cách.
\dapso{\pa{a} $6$ cách. \pa{b} $2$ cách.}}
